\documentclass[11pt,a4paper]{article}

\usepackage{fontspec}
\usepackage{xeCJK}
\setCJKmainfont{Hiragino Mincho ProN}
\setCJKsansfont{Hiragino Kaku Gothic ProN}
\usepackage[margin=2cm]{geometry}
\usepackage{amsmath,amssymb,amsfonts}
\usepackage{graphicx}
\usepackage{hyperref}
\usepackage{bm}
\usepackage{booktabs}
\usepackage{amsthm}

\newtheorem{theorem}{定理}
\newtheorem{proposition}[theorem]{命題}
\newtheorem{corollary}[theorem]{系}
\newtheorem{conjecture}[theorem]{予想}
\newtheorem{lemma}[theorem]{補題}

\newcommand{\Spec}{\mathrm{Spec}}
\newcommand{\Tr}{\mathrm{Tr}}
\newcommand{\Br}{\mathrm{Br}}
\newcommand{\Gal}{\mathrm{Gal}}
\newcommand{\Zp}{\mathbb{Z}[1/p]}
\newcommand{\Qp}{\mathbb{Q}_p}
\newcommand{\Zhat}{\hat{\mathbb{Z}}}
\newcommand{\Fp}{\mathbb{F}_p}
\newcommand{\Gm}{\mathbb{G}_m}

\begin{document}


\title{$\Spec(\mathbb{Z})$上の真空エネルギー条件の\\
  層理論的再定式化}

\author{Wright Brothers}
\affiliation{Independent Research}
\date{\today}

\begin{abstract}
本論文では、ゼータ正則化された真空エネルギーおよび弱エネルギー条件
(weak energy condition, WEC)を、算術的スキーム$\Spec(\mathbb{Z})$上の
層（sheaf）とコホモロジーの言語で再定式化する。
$\Spec(\mathbb{Z})$上に「エネルギー層」(energy sheaf) $\mathcal{F}$を
構成し、その大域切断がリーマンゼータ関数をエンコードし、
開部分スキーム$\Spec(\Zp)$への制限が$p$-除去ゼータ関数
$\zeta_{\neg p}(s) = \zeta(s)(1-p^{-s})$をエンコードすることを示す。
エタールコホモロジー (\'etale cohomology) の局所化完全列 (localization exact sequence) を用いて、
局所コホモロジー$H^2_{(p)}(\Spec(\mathbb{Z}), \Gm) \cong \Qp/\mathbb{Z}_p$
が「$p$成分のすべてのコホモロジー的障害」を担い、
局所化によりこの成分が除去されることを示す。
ボレル調整子 (Borel regulator) $\mathrm{reg}: K_{2n}(\mathbb{Z}) \to \mathbb{R}$
は$\zeta(1-2n)$に写像されるが、すべての素数$p$およびすべての$n \geq 1$に対して
局所化$\mathbb{Z} \to \Zp$の下で符号が反転することを証明する。
この符号反転が弱エネルギー条件の破れに対応すると予想し、
WECは連続的なエネルギー閾値ではなく、
「調整子の向き」(regulator orientation)---離散的な位相不変量---であると
提唱する。
これが正しければ、関連論文~\cite{paper_A,paper_C}で特定された
スケーリング問題は解消される：エネルギーを蓄積する必要はなく、
向きを反転させればよいのである。
局所化$\mathbb{Z} \to \Zp$の五つの物理的解釈
（ゲージ化、カルツァ--クライン、相転移、トポス論理変更、
およびアデール成分除去）を特定し、
トポロジカル材料におけるK理論的相転移が最も有望な
実験的実現であると提案する。
\end{abstract}

\maketitle

% ============================================================================
\section{序論}
\label{sec:intro}
% ============================================================================

関連論文~\cite{paper_A}において、「素数チャネル除去」(prime channel muting)
---素数$p$で割り切れるモードを射影除去する操作---が、
ゼータ正則化された真空エネルギーの符号を反転させることを示した：
すべての$p \geq 2$に対して$\zeta_{\neg p}(-3) < 0$である。
\cite{paper_C}では、この結果をボスト--コンヌ系 (Bost-Connes system) から
ワープドライブに至る十段階の理論的連鎖の中に位置づけ、
約$10^{69}$桁のスケーリング問題が未解決であることを特定した。

本論文では、同じ物理を$\Spec(\mathbb{Z})$上の代数幾何学と
層コホモロジーの言語で再定式化する。
動機は二つある：
\begin{enumerate}
  \item オイラー積$\zeta(s) = \prod_p (1-p^{-s})^{-1}$は
    単なる公式ではなく、\emph{スキーム}$\Spec(\mathbb{Z})$のゼータ関数であり、
    素数除去は\emph{局所化} (localization) という幾何学的操作に対応する。
  \item スケーリング問題は質的な解決を許容しうる：
    もしWECがエネルギー閾値（連続的）ではなく
    コホモロジー的障害（オン／オフ）であるならば、
    「反転」に比例するエネルギーは不要である。
\end{enumerate}

主張(2)は\emph{予想}であり、定理ではないことを強調する。
本論文はこの予想のための数学的枠組みを構築し、
それを確立するために何を証明する必要があるかを特定する。

% ============================================================================
\section{$\Spec(\mathbb{Z})$上のエネルギー層}
\label{sec:sheaf}
% ============================================================================

\subsection{構成}

$\Spec(\mathbb{Z})$上のザリスキ位相は、
$n \in \mathbb{Z} \setminus \{0\}$に対する
区別開集合$D(n) = \Spec(\mathbb{Z}[1/n])$を基底として持つ。
以下を定義する：

\begin{proposition}[エネルギー前層]
\label{prop:presheaf}
$\mathrm{Re}(s) > 1$なる$s \in \mathbb{C}$を固定する。
割り当て
\begin{equation}
  \mathcal{F}(D(n)) = \zeta(s) \cdot \prod_{p | n} (1 - p^{-s})
  \label{eq:energy_sheaf}
\end{equation}
および制限写像
$\mathrm{res}_{D(m) \to D(n)}: f \mapsto f \cdot \prod_{p|n,\, p \nmid m} (1-p^{-s})$
（$D(n) \subseteq D(m)$のとき）は、
$\Spec(\mathbb{Z})$上の$\mathbb{C}$値関数の前層を定義する。
\end{proposition}

\begin{proof}
前層の公理を検証する。(i) $\mathrm{res}_{U \to U} = \mathrm{id}$：
$m = n$のとき$\{p | n,\, p \nmid m\}$上の積は空である。
(ii) 推移性：$D(\ell) \subseteq D(n) \subseteq D(m)$に対して、
素因数が互いに素な集合に分解するため、因子は正しく合成される。
\end{proof}

前層$\mathcal{F}$は実際には\emph{層}（sheaf）である：
互いに素な素数集合上のオイラー積が因数分解するため、
貼り合わせ条件が成立する。

$\mathcal{F}(\Spec(\mathbb{Z})) = \zeta(s)$（大域切断）
および$\mathcal{F}(D(p)) = \zeta_{\neg p}(s)$（閉点$(p)$の
補集合上の切断）であることに注意する。

\subsection{解析接続}

$\mathrm{Re}(s) > 1$に対して$\mathcal{F}$は明確に定義される。
$\zeta(s)$の解析接続により、$\mathcal{F}$を
$s \in \mathbb{C}$をパラメータとする
$\Spec(\mathbb{Z})$上の有理型関数の層に拡張する。
負の奇数整数$s = 1 - 2n$における値が、
物理に関連する真空エネルギー値である。

% ============================================================================
\section{局所化完全列}
\label{sec:localization}
% ============================================================================

基本的な道具はエタールコホモロジーの局所化完全列である。
$i: \{(p)\} \hookrightarrow \Spec(\mathbb{Z})$
および$j: D(p) \hookrightarrow \Spec(\mathbb{Z})$を
それぞれ閉包含および開包含とする。
任意のエタール層$\mathcal{G}$に対して：
\begin{equation}
  \cdots \to H^n_{(p)} \to H^n(\Spec(\mathbb{Z}), \mathcal{G})
  \to H^n(D(p), \mathcal{G}) \to H^{n+1}_{(p)} \to \cdots
  \label{eq:les}
\end{equation}
ここで$H^n_{(p)} = H^n_{(p)}(\Spec(\mathbb{Z}), \mathcal{G})$
は$(p)$に台を持つ局所コホモロジーである。

$\mathcal{G} = \Gm$（乗法群層）に対して：

\begin{proposition}
\label{prop:local_cohom}
\begin{align}
  H^0(\Spec(\mathbb{Z}), \Gm) &= \mathbb{Z}^* = \{\pm 1\}, \\
  H^0(D(p), \Gm) &= \Zp^* = \{\pm 1\} \times \langle p \rangle, \\
  H^1(\Spec(\mathbb{Z}), \Gm) &= \mathrm{Pic}(\mathbb{Z}) = 0, \\
  H^2_{(p)}(\Spec(\mathbb{Z}), \Gm) &\cong \Qp/\mathbb{Z}_p.
  \label{eq:local_cohom}
\end{align}
\end{proposition}

$n=2$の段階における完全列は次を与える：
\begin{equation}
  0 \to \Qp/\mathbb{Z}_p \to \Br(\mathbb{Z})
  \to \Br(\Zp) \to 0.
  \label{eq:brauer_ses}
\end{equation}

\textbf{物理的解釈。}
$\Qp/\mathbb{Z}_p$は「$p$-局所的障害」の群である。
局所化$\mathbb{Z} \to \Zp$はこの群全体を
ブラウアー群 (Brauer group) から除去する。
$p$成分が非自明な任意の障害は、$p$の除去により弱められるか破壊される。

% ============================================================================
\section{K理論と調整子}
\label{sec:ktheory}
% ============================================================================

\subsection{$\mathbb{Z}$と$\Zp$のK群}

$\mathbb{Z}$の代数的K群（Quillen~\cite{Quillen1973}）は：
\begin{align}
  K_0(\mathbb{Z}) &= \mathbb{Z}, \quad K_1(\mathbb{Z}) = \mathbb{Z}/2, \\
  K_2(\mathbb{Z}) &= \mathbb{Z}/2 \quad \text{(Milnor~\cite{Milnor1971})}。
\end{align}

局所化の下で、$K_1$は自由部分を獲得する：
\begin{equation}
  K_1(\Zp) = \mathbb{Z}/2 \times \mathbb{Z},
  \label{eq:k1_local}
\end{equation}
ここで新しい$\mathbb{Z}$因子は$p \in \Zp^*$により生成される。
これはK群の\emph{位相的型}の変化であり：
除去された各素数に対して自由ランクが1増加する。

\subsection{ボレル調整子とゼータ値}

ボレル調整子~\cite{Borel1977}は準同型写像
\begin{equation}
  \mathrm{reg}: K_{2n-1}(\mathbb{Z}) \to \mathbb{R}^{d_n}
\end{equation}
であり、その像は格子をなし、その余体積は
$|\zeta(n)|$に（有理因子と$\pi$の冪を除いて）関連する。

本研究における鍵となる関係は：
\begin{equation}
  \zeta(1-2n) \sim \frac{|K_{2n-2}(\mathbb{Z})|}{|K_{2n-1}(\mathbb{Z})|} \cdot (\text{$\pi$の冪}),
  \label{eq:lichtenbaum}
\end{equation}
Lichtenbaum~\cite{Lichtenbaum1973}により確立され、
Voevodsky~\cite{Voevodsky2003}により（2-進の場合に）証明された。

\subsection{調整子の符号反転}

\begin{theorem}[調整子の符号反転]
\label{thm:reg_flip}
任意の素数$p \geq 2$および任意の正整数$n$に対して、
「局所化された調整子値」
$\zeta_{\neg p}(1-2n) = \zeta(1-2n) \cdot (1 - p^{2n-1})$
は$\zeta(1-2n)$と反対の符号を持つ。
\end{theorem}

\begin{proof}
$p \geq 2$, $n \geq 1$のとき$1 - p^{2n-1} < 0$である。
\end{proof}

数値的な内容は\cite{paper_A}の符号反転定理と同一であるが、
K理論的解釈は新しい：

\begin{conjecture}[調整子の向き]
\label{conj:reg}
$\zeta(1-2n)$の符号は位相不変量である：
$\mathbb{R}^{d_n}$におけるボレル調整子格子の\emph{向き}（orientation）である。
局所化$\mathbb{Z} \to \Zp$はこの向きを反転させる。
物理的には、この向きが真空エネルギーの正（WEC充足）または
負（WEC破れ）を決定する。
\end{conjecture}

予想~\ref{conj:reg}が正しければ、WECは算術的スキームの
離散的（オン／オフ）不変量であり、連続的なエネルギー閾値ではない。
WECを「破る」とは調整子の向きを反転させることを意味し、
それは単一の局所化操作により達成される。

% ============================================================================
\section{局所化の物理的解釈}
\label{sec:physics}
% ============================================================================

操作$\mathbb{Z} \to \Zp$は、数学的に等価な五つの
異なる物理的解釈を許容する：

\textbf{(1) $\mathbb{Z}/p\mathbb{Z}$対称性のゲージ化。}
$\Zp$において$p$は単元であり、したがって$|n\rangle$と$|pn\rangle$は
ゲージ同値となる。独立な自由度はゲージ軌道
$\{n \cdot p^k : k \in \mathbb{Z}\}$であり、
$p$と互いに素な整数でラベルされる。

\textbf{(2) カルツァ--クライン非コンパクト化。}
各素数$p$は有効半径$R_p \propto 1/\log(p)$の
コンパクト余剰次元に対応する
（BCハミルトニアン$H|n\rangle = \log(n)|n\rangle$より）。
$p$での局所化はこの次元を非コンパクト化し（$R_p \to \infty$）、
エネルギースペクトルに対するKK制約を解放する。

\textbf{(3) 相転移。}
オイラー因子$(1-p^{-\beta})^{-1}$は素数$p$に対する秩序パラメータである。
局所化は$\beta_p \to \infty$（選択的凍結）に対応し、
部分的な相転移である。

\textbf{(4) トポス論理変更。}
$\Spec(\mathbb{Z})$と$\Spec(\Zp)$は異なる
エタールトポスを定義し、異なる内部論理（ハイティング代数）を持つ。
命題（WECなど）は異なるトポスにおいて異なる真理値を持ちうる。

\textbf{(5) アデール成分除去。}
アデール環$\mathbb{A}_\mathbb{Q} = \mathbb{R} \times \prod'_p \Qp$。
局所化は$\Qp$成分を除去する：真空はその量子揺らぎの
「$p$進チャネル」を失う。

五つの解釈はすべて同じ数学的操作を記述する。
実験的に最もアクセスしやすいのは\textbf{(1)}：
離散ゲージ理論は超伝導量子ビット配列で実現可能であること、
および\textbf{(3)}：BCハミルトニアンの量子シミュレータで
相転移を誘起できることである。

% ============================================================================
\section{境界：壁の物理}
\label{sec:boundary}
% ============================================================================

局所化が空間の\emph{有限領域}に適用される場合、
$\Spec(\mathbb{Z})$（外部、WEC = 真）と
$\Spec(\Zp)$（内部、WEC = 偽）の間の境界は
閉点$V(p) = \Spec(\Fp)$である。

\subsection{壁のデータとしての局所コホモロジー}

局所コホモロジー$H^2_{(p)} \cong \Qp/\mathbb{Z}_p$
（命題~\ref{prop:local_cohom}）は壁の上に台を持つ。
$\Qp/\mathbb{Z}_p$は階層的構造を持つ無限群である：
\begin{equation}
  \Qp/\mathbb{Z}_p = \bigcup_{n=1}^{\infty} p^{-n}\mathbb{Z}_p / \mathbb{Z}_p
  \cong \bigcup_{n=1}^{\infty} \mathbb{Z}/p^n\mathbb{Z}.
\end{equation}

各レベル$n$は$p^n$個の「壁モード」を担う。
これらは\emph{算術的エッジ状態}(arithmetic edge states)：
二つの算術的相の界面に局在する境界自由度である。

\subsection{トポロジカルエッジ状態との類似}

トポロジカル絶縁体において、異なるK理論不変量を持つ
二つの相の間の境界はギャップレスエッジ状態を支える
（バルク--境界対応~\cite{Kitaev2009}）。

ここで、$K_1(\mathbb{Z}) = \mathbb{Z}/2$および
$K_1(\Zp) = \mathbb{Z}/2 \times \mathbb{Z}$：
K群の差は$\Delta K_1 = \mathbb{Z}$である。
バルク--境界対応により、壁は無限のエッジ状態の塔を
支えるはずである---これはまさに$\Qp/\mathbb{Z}_p$により
数え上げられるモードである。

\subsection{壁における予測現象}

\begin{enumerate}
  \item \textbf{算術的エッジ状態}：$\Qp/\mathbb{Z}_p$により
    数え上げられる境界局在モード。
  \item \textbf{輻射}：壁における対生成（ホーキング輻射に類似）、
    温度$T_{\mathrm{wall}} \sim \hbar\omega_0 \log(p) / (2\pi k_B)$。
  \item \textbf{屈折率不連続}：壁を横切るとモード密度が
    $\sqrt{1-1/p}$倍変化し、
    有効屈折率比$n_{\mathrm{in}}/n_{\mathrm{out}} = \sqrt{1-1/p}$を与える。
  \item \textbf{安定性}：壁はトポロジカルに保護される
    （異なるK理論相は連続的に変形できない）が、
    標数$p$における暴分岐を支え、
    最小壁厚に制限を課す。
\end{enumerate}

これらの予測は数学的枠組みの帰結であり、
\cite{paper_A}の算術的真空工学プログラムが
レベル3以上で成功した場合に検証可能である。

% ============================================================================
\section{議論}
\label{sec:discussion}
% ============================================================================

\subsection{層の再定式化は予測能力を向上させるか？}

本論文の数値的予測は\cite{paper_A}のものと同一である：
$\zeta_{\neg p}(-3) = (1-p^3)/120$。
しかし、層理論的枠組みは以下を提供する：
\begin{enumerate}
  \item 符号反転の\emph{質的}解釈としての
    調整子の向きの反転（予想~\ref{conj:reg}）。
    これが正しければスケーリング問題は解消される。
  \item 新しい実験的予測：K理論的相転移
    $K_1(\mathbb{Z}) \to K_1(\Zp)$は、
    凝縮系アナログにおけるトポロジカル相転移として
    観測可能なはずである。
  \item 壁の物理の枠組み：局所コホモロジー
    $\Qp/\mathbb{Z}_p$は具体的な境界現象を予測する。
\end{enumerate}

\subsection{未解決の予想}

\begin{enumerate}
  \item 予想~\ref{conj:reg}：WECと調整子の向きの同一視は
    証明されていない。
  \item 壁の現象（第~\ref{sec:boundary}節）は、
    数学的なバルク--境界対応が物理的な対応物を持つことを仮定している。
  \item スケーリング問題が「解消される」という主張は、
    完全に予想~\ref{conj:reg}に依存する。
    もしWECが離散的（位相的）条件ではなく
    本質的に連続的（エネルギー的）条件であるならば、
    スケーリング問題は残る。
\end{enumerate}

% ============================================================================
\section{結論}
\label{sec:conclusion}
% ============================================================================

$\Spec(\mathbb{Z})$上にリーマンゼータ関数のオイラー積構造を
エンコードするエネルギー層$\mathcal{F}$を構成し、
エタールコホモロジーを用いてその局所化特性を計算し、
すべての素数$p$に対して
$\mathbb{Z} \to \Zp$の下でボレル調整子が符号反転を
起こすことを示した。

この符号反転は離散的な位相不変量---調整子格子の向き---であり、
弱エネルギー条件の充足または破れに対応すると予想する。
これが正しければ、WECは$\Spec(\mathbb{Z})$の算術幾何学的性質として
再構成され、所望の効果に比例するエネルギーを必要とせず
局所化により「切り替え」可能となる。

局所化の物理的実現は五つの等価な解釈を許容し、
そのうちK理論的トポロジカル相転移と離散ゲージ理論シミュレーションが
実験的に最もアクセスしやすい。
二つの算術的相の境界は局所コホモロジー$\Qp/\mathbb{Z}_p$により
特徴づけられ、具体的なエッジ状態現象を予測する。

この枠組みはスケーリング問題を、定量的課題
（$10^{69}$桁のエネルギーの蓄積）から定性的課題
（調整子の向きを反転させるトポロジカル相転移の実現）へと
変換する。この定性的解決が物理的に妥当であるか否かが、
中心的な未解決問題として残る。

\section*{謝辞}
本研究はWright Brothers共同研究グループにより
独立研究として遂行された。

\begin{thebibliography}{30}

\bibitem{paper_A}
Wright Brothers,
``Arithmetic vacuum engineering: Prime channel muting via
Bost-Connes phase transitions,''
companion paper (2026).

\bibitem{paper_C}
Wright Brothers,
``From Bost-Connes to warp drive: A theoretical chain,''
companion paper (2026).

\bibitem{Quillen1973}
D.~Quillen,
in \emph{Algebraic $K$-theory I}, Springer LNM \textbf{341},
pp.~85--147 (1973).

\bibitem{Milnor1971}
J.~Milnor,
\emph{Introduction to Algebraic $K$-theory},
Princeton University Press (1971).

\bibitem{Borel1977}
A.~Borel,
``Cohomologie de $\mathrm{SL}_n$ et valeurs de fonctions z\^eta
aux points entiers,''
Ann. Scuola Norm. Sup. Pisa \textbf{4}, 613 (1977).

\bibitem{Lichtenbaum1973}
S.~Lichtenbaum,
``Values of zeta-functions, \'etale cohomology, and algebraic
$K$-theory,''
in \emph{Algebraic $K$-theory II}, Springer LNM \textbf{342},
pp.~489--501 (1973).

\bibitem{Voevodsky2003}
V.~Voevodsky,
``Motivic cohomology with $\mathbb{Z}/2$-coefficients,''
Publ. Math. IHES \textbf{98}, 59 (2003).

\bibitem{Kitaev2009}
A.~Kitaev,
AIP Conf. Proc. \textbf{1134}, 22 (2009).

\end{thebibliography}

\end{document}
